\begin{enumerate}[label=\thesubsection.\arabic*,ref=\thesubsection.\theenumi]
\item The statements (S1), (S2), and (S3) pertain to the scores obtained by students in
an exam. The maximum possible marks in the exam is 150.
	\begin{enumerate}
		\item[(S1)] The highest score is 100.
		\item[(S2)] The fourth highest score is 76.
		\item[(S3)] There are at least four students whose scores are within 25 of each other.
	\end{enumerate}
Which of the following options are necessarily correct?
	
	\begin{enumerate}
		\item (S1) and (S2) together imply (S3)
		\item (S1) and (S3) together imply (S2)
		\item (S2) and (S3) together imply (S1)
		\item (S1) implies (S3)
	\end{enumerate}
	





%--------------------------Q2------------------------------------------------------------

\item Real numbers $y$, $p$, and $n$ (all greater than 1) satisfy
	$(log_{p^{\frac{1}{n}}} y)(log_{y^{\frac{1}{n}}} p) = 16$,
		where the logarithms are taken to the bases $p^{\frac{1}{n}}$ and $y^{\frac{1}{n}}$ .
The value of $n$ is
		\begin{multicols}{4}
			\begin{enumerate}
				\item $2$
				\item $4$
				\item $8$
				\item $16$
			\end{enumerate}
		\end{multicols}



%================================Q3=====================================================

\item Three children P, Q, R and two grown-ups X, Y play a badminton doubles
tournament. X and Y are parents to two of the children playing. The child of X is
not the same as the child of Y. Exactly one of the children does not have a parent
playing in the tournament. The following rules are followed:
	\begin{enumerate}[label=\roman*]
		\item A parent and his/her child cannot be on the same team.
		\item A match can feature at most one parent and his/her child, that is, a
maximum of one parent-child pair can play in a match.
	\end{enumerate}
The following matches were played:\\
	\begin{tabular}{|l|l|l|}
		\hline
		&TEAM 1&TEAM 2\\\hline
		MATCH 1&P and X&Q and R\\\hline
		MATCH 2&P and R&X and y\\\hline
		MATCH 3&R and X&Q and Y\\\hline
	\end{tabular}\\
\par Which of the following options is correct?

\begin{enumerate}
	\item P does not have any parent playing
	\item Q does not have any parent playing
	\item R does not have any parent playing
	\item X does not have a child playing
\end{enumerate}
			


%===========================Q4=============================
\item Let $P_k$ represent the perimeter of a square with sides of length $k$. The value of the
	expression $P_1 + P_2 + P_3 + ...+ P_{10}$ is:
	
	\begin{multicols}{4}
	\begin{enumerate}
		\item 55
		\item 110
		\item 220
		\item 440
	\end{enumerate}
	\end{multicols}





%================================Q5=======================================================
\item Consider the differential equation $\dot{\vec {w}} = A\vec{w}$ , with$\vec{w}(t = 0)=\myvec{1\\1}$
.If
$w(t)=e^t\vec{u_x}+e^{-2t}\vec{u_y}$ be the solution to the equation where$\vec{u_x}$ and$\vec{u_y}$ are
unit vectors along the positive x and y axes respectively, then which of the
following options is the correct matrix representing A?
\begin{multicols}{4}
\begin{enumerate}
	\item $\myvec{1&0\\
		0&-2}$
	\item $\myvec{-1 &0\\
		0&-2}$
	\item $\myvec{0&-2\\
		1&0}$
	\item $\myvec{0&2\\
		-1&0}$
\end{enumerate}
\end{multicols}
%==================================Q6=============================================================
\item A surface is given by $z^2=2x^2-y^2$ and$\vec{n}$ and $\vec{-n}$ are the unit vectors
are unit normal vectors to the surface at the point $\vec{P}=\vec{i}+\sqrt{2}\vec{k}$
Which of the following vectors can be$\vec{n}$, where$\vec{i},\vec{j}$ and $\vec{k}$
along x, y and z axes, respectively?
\begin{multicols}{4}
	\begin{enumerate}
		\item $\vec{i}-\sqrt{2}\vec{k}$
		\item $\frac{2}{3}\vec{i}-\frac{1}{3}\vec{k}$
		\item $\sqrt{2}\vec{i}-\sqrt{3}\vec{k}$
		\item $\frac{\sqrt{2}\vec{i}-\vec{k}}{\sqrt{3}}$
	\end{enumerate}
\end{multicols}
%========================================Q7====================================================
\item The Laplace Transform of the signal$ x\brak{t}=u\brak{t-2}*\brak{tu\brak{t}}$ is given by
which of the following expressions?\\
\{"*" represents the convolution operator\}
\begin{multicols}{4}
	\begin{enumerate}
		\item $\frac{e^{-2s}}{s^2\brak{s-2}}$
		\item $\frac{e^{-2\brak{s-2}}}{s^3}$
		\item $\frac{se^{-2s}}{\brak{s-2}^2}$
		\item $ \frac{e^{-2s}}{s^3}$
	\end{enumerate}
\end{multicols}
%==========================================================Q8=====================================
\item Two analog signals $x_1\brak{t}$and $x_2\brak{t}$ (t in second), are sampled at a rate
$F_s = 40$ Hz, where
$x_1 \brak{t} = \cos \brak{20\pi t},t\geq0,$ and $x_2 \brak{t} = \cos \brak{100\pi t}, t\geq0$.
The first ten samples (starting from t = 0) are considered for the analysis.
Which of the following statements is TRUE? 
	\begin{enumerate}
		\item All of the first three samples of $x_1\brak{t}$ are greater than the corresponding samples of $x_2 \brak{t}$. 
		\item All of the last three samples of $x_1\brak{t}$ are greater than the corresponding samples of $x_2 \brak{t}$. 
		\item All of the samples of $x_2\brak{t}$ are greater than the corresponding samples of $x_1 \brak{t}$.
		\item All of the fourth to seventh samples of $x_1\brak{t}$ are equal to the corresponding samples of $x_2 \brak{t}$.
	\end{enumerate}
%===========================================================Q9=========================================
\item The response of a discrete time system $y[n]$ obeys the following relation:\\
$y[n] =\frac{5}{6} y[n-1]-\frac{1}{6}y[n-2] +x[n]$.\\
The input to the system is $x[n] = \delta[n] - \frac{1}{3}\delta[n-1]$.

Which of the following options is TRUE for $y[n]$?
	\begin{enumerate}
		\item Stable and causal response
		\item Stable and non-causal response
		\item Unstable and causal response
		\item Unstable and non-causal response
	\end{enumerate}
%========================================================Q10==========================================
\item A control system is shown in the Figure.
	Which option represents the correct transfer function of the system?
	%DRAW DIAGRAM!!!!!!!!!!!!!!!!!!!!!!!!!!!!!!
	\begin{multicols}{4}
		\begin{enumerate}
			\item $\frac{1}{\brak{s+4}^2}$
			\item $\frac{1}{s+4}$
			\item $\frac{2}{s+4}$
			\item $\frac{1}{\brak{s^2+8s+17}}$
		\end{enumerate}
	\end{multicols}
%==================================================Q11================================================
\item Consider a discrete memoryless source with an alphabet of four source symbols.
	$s\brak{t}$ is a multi-level $\brak{-1, 0, +1, +2}$ signal representing a long sequence of random
symbols from the above source which is generating 104 symbols per second.
Which of the following options is the correct value of equivalent Nyquist
bandwidth of $s\brak{t}$?
\begin{multicols}{4}
	\begin{enumerate}
		\item 10kHz
		\item 64kHz
		\item 5kHz
		\item 20kHz
	\end{enumerate}
\end{multicols}
%=======================================================Q12===========================================
\item Consider the matrix $\vec{M}=\myvec{2&1&1\\1&3&0\\-1&a&b}$.Which of the following options is/are TRUE if det$\brak{\vec{M}}\neq0$?
\begin{multicols}{4}
	\begin{enumerate}
		\item a=$\frac{-1}{2}$and b=$\frac{-1}{2}$
		\item a=$\frac{1}{2}$and b=$\frac{1}{2}$
		\item a=$-3$and b=$0$
		\item a=$\frac{1}{2}$and b=$-3$
	\end{enumerate}
\end{multicols}

%=======================================================Q13============================================
\item Consider the circuit shown in the Figure, where the input $v_i \brak{t}$ is in Volt.
The average power (in mW) dissipated in the load resistance of $1 k\Omega$ at the
resonant frequency is:
(rounded off to two decimal places)
%DRAW DIAGRAM 
%=========================================================Q14=======================================
\item Consider the two series, SA and SB, where
	\begin{align*}
		S_A&=\sum_{n=1}^\infty \frac{n^2}{2^n}\\
		S_B&=1+\frac{1}{2}+\frac{1}{8}+\frac{1}{16}+\frac{1}{64}+\frac{1}{128}+\frac{1}{512}+....
	\end{align*}
Which of the following statements is correct for the two given series?
\begin{enumerate}
	\item Both $S_A$ and $S_B$ converge.
	\item Neither $S_A$ nor $S_B$ converges.
	\item $S_A$ converges but $S_B$ does not converge.
	\item $S_B$ converges but $S_A$ does not converge.
\end{enumerate}
%=========================================================Q15===============================================
\item The continuous time signal $x\brak{t}$ is real, periodic with period T and satisfies the
Dirichlet conditions.
The Fourier series representation of $x\brak{t} = \sum_{-\infty}^\infty a_ne^{j\brak{\frac{2\pi n t}{T}}}$
and $x\brak{t}$ satisfies the following:$x\brak{t=\frac{T}{2}}=-x\brak{t}$. For any integer m, which of the following options is correct? 
\begin{multicols}{4}
	\begin{enumerate}
		\item $a_{2m}=0$
		\item $a_{2m}=1$
		\item $a_{2m}=a_{2m+1}$
		\item $a_{2m}=-1$
	\end{enumerate}
\end{multicols}
%==========================================================Q16=============================================
\item Let X, N, Y and Z be random variables. The variables X and N are independent of
each other. X is uniformly distributed between -1 and 1; N follows Normal
distribution with zero mean and unity variance.\\
Y and Z are defined as,$Y = X+N$ and $Z = X^2 +N$.\\
Which of the following pairs represents the values of correlation between X and Y
and that between X and Z? 
\begin{enumerate}
	\item $\frac{1}{3}$ and 0
	\item $\frac{1}{3}$ and $\frac{1}{9}$
	\item $\frac{1}{3}$ and $\frac{1}{3}$
	\item 1 and 0
\end{enumerate}
%==================================================Q17=====================================================
\item In the given circuit, $L = 1 \mu H$ and $C = 1 \mu F$. The phasor diagram for IC and IL is
	also shown. Assume that the phase $\brak{\theta_1 + \theta_2}$ is $90^\circ$ at a frequency of 159.15 kHz.
Among the following options, what is the nearest integer value of$ R_C\times  R_L$?
\begin{multicols}{4}
	\begin{enumerate}
		\item 0
		\item 1
		\item 2
		\item 10
	\end{enumerate}
\end{multicols}
%======================================================Q18================================================
\item For the control system shown in the Figure, the transfer function of a plant,

	$G\brak{s} = \brak{s+1}\brak{s+2}$ is connected in cascade with a compensator
$c\brak{s} = K \brak{s + \alpha}$, where $K $and$ \alpha $are positive real valued constants.
Which of the following pairs $\brak{K, \alpha}$ represent the correct values for the closed
loop system to have poles at $\brak{-3 \pm j\sqrt{5}}$?
%NEED TO DRAW DIAGRAM
\begin{multicols}{4}
	\begin{enumerate}
		\item 2,3
		\item 3,4
		\item 2,4
		\item 3,3
	\end{enumerate}
\end{multicols}


%===============================================Q19====================================================
\item The state and output equations for a control system are:
	\begin{align*}
		\dot{x}&=\myvec{-4&-1.5\\4&0}x+\myvec{2\\0}u\\
		y=\myvec{1.5 & 0.625}x
	\end{align*}
Which of the following expressions correctly represents the transfer function$\frac{Y\brak{s}}{u\brak{s}}$
of the system with zero initial conditions?
\begin{multicols}{4}
	\begin{enumerate}
		\item $\frac{3s}{s^2+4s-6}$
		\item $\frac{3s+5}{s^2+4s-6}$
		\item $\frac{3s+5}{s^2+4s+6}$
		\item $\frac{3s}{s^2+4s+6}$
	\end{enumerate}
\end{multicols}
%================================================================Q20=====================================
\item The address of the first location of a 256 kilo byte (KB) memory is $\brak{2500}_H$.
Choose the correct address of the last location of the memory.
\begin{multicols}{4}
	\begin{enumerate}
		\item $\brak{2FFF}_H$
		\item $\brak{124FF}_H$
		\item $\brak{424FF}_H$
		\item $\brak{324FF}_H$
	\end{enumerate}
\end{multicols}
%======================================================================Q21=================================
\item Consider a real signal $x\brak{t}, -\infty < t < \infty$, such that
	$x\brak{t} = 0$ for$t <0, x\brak{t} = 2$ for $0 \leq t < 1$ and $x\brak{t} = 0$ for $t \geq 1$.

Let$ E[x\brak{t}] = \int_{-\infty}^{\infty}x\brak{t}]2 dt$.
Which of the following options correctly represents the ratio,
$E[x\brak{t}]⁄E[3x\brak{-3t + 5}]$?
\begin{multicols}{4}
	\begin{enumerate}
		\item 3
		\item 1
		\item $\frac{1}{3}$
		\item $\frac{1}{9}$
	\end{enumerate}
\end{multicols}
%========================================================Q22==============================================
\item A QPSK modulated signal from an additive white Gaussian noise (AWGN) channel is received with an $E_b/N_o = 8.4$ dB at the input of a coherent QPSK demodulator. A maximum-likelihood reception method is used in the
demodulator.
Assume the complimentary error function
\begin{align*}
	e rfc\brak{u} \approxeq [1⁄\brak{u\sqrt{\pi}}]exp\brak{-u^2}
\end{align*}.
Which is the nearest bit error rate (BER) at the output of the demodulator?
\begin{multicols}{4}
	\begin{enumerate}
		\item $10^{-3}$
		\item $10^{-4}$
		\item $10^{-5}$
		\item $10^{-6}$
	\end{enumerate}
\end{multicols}
%==========================================Q23=============================================================
\item What is the 10's complement of $\brak{47}_10$?
	\begin{multicols}{4}
		\begin{enumerate}
			\item 52
			\item 53
			\item 54
			\item 55
		\end{enumerate}
	\end{multicols}
%=======================================Q24===============================================================
\item Consider a real baseband signal$x\brak{t} =e^{-2t}$ , for t (in seconds) $\geq$ 0.
	If 99\% of energy of $x\brak{t}$ lies within B Hz, then which of the following options is
TRUE for the value of B? 
\begin{enumerate}
	\item B> 1kHz
	\item $63/\pi$Hz<B<$64/\pi$Hz
	\item $126/\pi$Hz<B<$128/\pi$Hz
	\item $B$<1Hz
\end{enumerate}
%=====================================================Q25=================================================
\item Consider the discrete time system (S) with input x[n] and output y[n] as shown in
	the Figure. The two sub-systems represented by their impulse responses $h_1[n]$
and $h_2[n]$ are linear and time invariant. Which of the following statements is necessarily TRUE?
%DRAW DIGRAM
\begin{enumerate}
	\item $S$ is causal.
	\item $S$ is linear and time invariant.
	\item $S$ is linear and time varying.
	\item $S$ is non-linear.
\end{enumerate}
%=======================================================Q26================================================
\item A Boolean function, $f\brak{x,y,z}$ with x as MSB and z as LSB is realized by 4:1
	multiplexer (MUX) with select lines, $S_1 and S_0$ ( $S_1$ is MSB, $S_0$ is LSB) and inputs,
$I_0, I_1, I_2, I_3$ as shown in the Figure.
Which of the following options is the correct expression of $f\brak{x,y,z}$?
%DRAW DIAGRAM 
\begin{multicols}{4}
	\begin{enumerate}
		\item $xz+y$
		\item $x\bar{y}+z$
		\item $xy+\bar{z}$
		\item $\bar{x}y+\bar{z}$
	\end{enumerate}
\end{multicols}
%============================================Q27===========================================================
\item A shift-left Shift Register (SR) and a D flip-flop are connected to a synchronized
clock as shown in the Figure. Assume that the SR and D flip-flops are initially
cleared and the XOR gate has no propagation delay.
Which of the following options gives the correct binary representation
$\brak{b_7b_6b_5b_4b_3b_2b_1b_0}$ of the content of the shift register immediately after the
5\textsuperscript{th} clock transition (positive edge)?
%DRAW DIAGRAM
\begin{multicols}{4}
	\begin{enumerate}
		\item 00011111
		\item 10111111
		\item 00111111
		\item 11000011
	\end{enumerate}
\end{multicols}
%==================================================Q28=====================================================
\item A small signal source,$ V_i \brak{t} = A \cos\brak{105t } + B\sin\brak{107 t}$ is applied to a
BJT circuit as shown in the Figure.
Assume zero source resistance, $V_{BE} = 0.7 $V, $\beta_{dc} = 99$, Early voltage = 100 V and
Thermal voltage = 25 mV. Effect of internal parasitic capacitances of the BJT
may be neglected.
Which expression is the best approximation of the output voltage $V_o\brak{t}$?
%DRAW DIAGRAM
\begin{enumerate}
	\item $-9.1[A\cos\brak{10^5t}+B\sin{10^7t}$
	\item $9.1[A\cos\brak{10^5t}-B\sin{10^7t}$
	\item $-190.4[A\cos\brak{10^5t}+B\sin{10^7t}$
	\item $190.4[A\cos\brak{10^5t}-B\sin{10^7t}$
\end{enumerate}
%==================================================Q29===================================================
\item Consider the four-variable Boolean function,
	$f\brak{w,x,y,z} = \sum m\brak{0,2,5,7,8,10,13,14,15}$ with ‘w’ as MSB and ‘z’ as LSB.
Which of the following expressions is/are the valid form(s) of$ f\brak{w, x, y, z}$?
\begin{enumerate}
	\item $\bar{x}\bar{z}+xz+wxy$
	\item $xz+wxy+w\bar{x}\bar{z}+\bar{w}xy\brak{z}$
	\item $\bar{x}\bar{z}+wxy+w\bar{x}\bar{z}+\bar{w}xy\bar{z}$
	\item $\bar{x}\bar{y}+xz+wy\bar{z}$
\end{enumerate}
%===============================================================Q30=======================================
\item Let $x_1 \brak{t} = \cos\brak{2\pi n t}$ and $x_2\brak{t} = 2\sin\brak{4\pi n t}$ represent two sinusoids
for a positive integer $n$ and $-\infty< t< \infty$ .
Which of the following statements about $x_1\brak{t}$ and $x_2\brak{t}$ is/are valid? 
\begin{enumerate}
	\item $x_1 \brak{t}$ and $x_2 \brak{t}$are orthogonal to each other over $0 \leq t < 1/n$ . 
	\item $x_1 \brak{t}$ and $x_2 \brak{t}$are orthonormaal to each other over $0 \leq t < 1/n$ .
	\item $x_1 \brak{t}$ is a harmonic of $x_2 \brak{t}$.
	\item $x_1 \brak{t}$ and $x_2 \brak{t}$ are non-orthogonal to each other over $0 \leq t < 1/\brak{2n}$ .
\end{enumerate}
%==============================================================Q31==========================================
\item Consider the unity negative feedback control system shown in the Figure.
The value of gain K (>0) at which the given system will remain marginally stable
is :
(Answer in integer)
%DRAW DIAGRAM
%=========================================================Q32===============================================
\item Consider the square region R in the X-Y plane as shown with the dark shading in
	the Figure. The value of$\int\int_{R} \brak{x_2 + y_2 - 1}dxdy$ is .
(rounded off to two decimal places) 
%DRAW DIAGRAM
%=======================================================Q33==================================================
\item Consider the ideal diodes D1 and D2 as shown in the Figure with cut-in voltage
	$V_\gamma = 0$ Volt and $v_i\brak{t}$ is in Volt.
		The maximum voltage (Volt) of the output $v_o\brak{t}$ is .
(rounded off to two decimal places)
%DRAW DIAGRAM
\end{enumerate}

